%\vspace*{1cm}
\begin{minipage}{\linewidth}
\selectlanguage{french} 
\begin{abstract}
Ce projet vise à concevoir et implémenter une plateforme PaaS innovante pour le déploiement en un clic d'applications web, en utilisant un lien GitHub public ou privé via une authentification OAuth sécurisée. Le déploiement s'effectue dans un environnement cloud orchestré par Kubernetes, qui gère l'orchestration des conteneurs Docker pour assurer scalabilité, haute disponibilité (HA) avec auto-scaling horizontal basé sur les métriques de charge, et une stratégie multi-tenancy adoptant l'approche silo pour isoler complètement les ressources des utilisateurs, garantissant ainsi une sécurité renforcée et une confidentialité des données. L'expérience utilisateur est centrée sur une interface web intuitive et responsive, permettant non seulement de soumettre le dépôt Git pour un build automatique, mais aussi de monitorer en temps réel les logs d'exécution, les métriques de performance (CPU, mémoire, trafic réseau), et de générer dynamiquement des sous-domaines avec activation automatique de certificats SSL via Let's Encrypt pour une accessibilité sécurisée. De plus, la plateforme intègre une gestion avancée des environnements, incluant la configuration de variables d'environnement (env) et de secrets sensibles (comme les clés API ou mots de passe) stockés de manière chiffrée dans Kubernetes Secrets, évitant toute exposition accidentelle. Un historique des déploiements est maintenu, offrant des fonctionnalités de rollback vers des versions précédentes en cas d'erreur, avec un journal détaillé des builds, tests et mises en production pour une traçabilité complète. La persistance des volumes est assurée via des PersistentVolumes dans Kubernetes, permettant aux applications stateful (comme celles avec bases de données) de conserver les données malgré les redéploiements ou scaling. Inspiré de solutions comme Railway, Render ou Vercel, ce système supporte divers langages et frameworks (Node.js, Python, etc.), intègre des pipelines CI/CD automatisés pour des mises à jour continues via webhooks GitHub, et favorise une adoption agile dans des contextes DevOps, réduisant les temps de mise en production tout en optimisant les coûts cloud via une allocation dynamique des ressources, rendant la plateforme adaptée à des environnements hybrides ou multi-cloud
\end{abstract}
\textbf{Mots Clée:} PaaS, Déploiement automatique, GitHub, Git, Cloud computing, Kubernetes, Haute disponibilité, SSL.
\end{minipage}